% ===================================================================
%  तक्ता क्रमांक १० — समुद्र. बंदर.
%
%  Traduction, faite en 2026, en marathi d'aujourd'hui, sur l'ido de
%  texto/io. Regles de la colonne : voir 10-takta-01.tex. Une seule
%  scene, vingt-cinq alineas numerotes : les cles ne portent pas
%  d'indice de scene.
%
%  L'AMIRAL S'APPELLE सरखेल, ET C'EST LE PLUS BEAU MOT DE LA
%  COLONNE. Ce n'est pas une traduction ni un emprunt : c'est le
%  titre que portait le chef de la flotte marathe, celui de Kanhoji
%  Angre, qui tint la cote de Konkan contre les compagnies
%  europeennes au debut du XVIIIe siecle. Le livret ecrit
%  « admiralo », l'internationalisme ; le marathi avait un mot a lui,
%  et il est plus ancien. De meme le quartier-maitre est तांडेल, le
%  chef de barque, et l'escadre est आरमार — le nom meme de la marine
%  marathe.
%
%  LA COTE ET LE PORT SONT ENTIEREMENT MARATHES, parce que le Konkan
%  est une cote de marins :
%
%    बंदर        le port          गोदी         le chantier, le dock
%    धक्का        le wharf         घाट          le quai
%    खाडी (—)    l'estuaire       गलबत (13)    le voilier
%    आगबोट (3)   le vapeur        होडी (46)    la barque
%    डोलकाठी (51) le mat           वल्हे (47)   la rame
%    शिड         la voile         सुकाणू (30)  le gouvernail
%    यारी (62)   la grue          खलाशी        le matelot
%    दीपगृह (93) le phare         लाट (72)     la vague
%    भरती / ओहोटी le flux et le reflux
%    बेट (99)    l'ile            पठार (96)    le plateau
%    बुरूज       le bastion       तट (15)      le rempart
%    शिबंदी      la garnison      तोफ (22)     le canon
%    समशेर       le sabre         दुर्बीण      la longue-vue
%
%  ET नांगर EST A LA FOIS L'ANCRE ET LA CHARRUE. Le meme mot que le
%  tableau 7 employait pour le labour sert ici pour les ancres (77)
%  qui ornent le balcon du cercle. Ce n'est pas une commodite : les
%  deux objets mordent le sol pour tenir, et le marathi les a nommes
%  d'un seul mot.
%
%  ET घाट AUSSI SERT DEUX FOIS : le defile de montagne au tableau 9,
%  le quai au bord de l'eau ici. Les deux sont marathes, les deux
%  sont courants, et aucun des deux n'est une metaphore de l'autre.
%
%  LE MATERIEL MODERNE, LUI, S'EMPRUNTE — croiseur, torpilleur,
%  fregate, tram, camion : la marine a vapeur n'a pas de nom marathe
%  et n'a pas a en recevoir un. Meme regle qu'au marron chaud du
%  tableau 7.
%
%  UNE INVERSION PREVENUE — le pont du voilier — ET SEPT QUI ONT
%  PASSE. C'est la troisieme fois de suite que le compte tombe sur
%  a peu pres le meme rapport, et c'est instructif : la postposition
%  ne se laisse pas anticiper, elle se laisse verifier.
%
%  ET UN CONTROLE A ETE CORRIGE PLUTOT QUE DESARME. kolonoj.py
%  signalait « दृश्य » comme du hindi ; c'est du marathi parfaitement
%  ordinaire — une vue, un spectacle — et ce tableau l'emploie deux
%  fois a bon droit. La regle etait trop large : ce qui serait
%  fautif est « पहिले दृश्य » a la place de « पहिला प्रवेश ». Elle vise
%  maintenant le mot PRECEDE DE SON ORDINAL. On ne desarme pas un
%  controle qui crie a tort, on le vise mieux.
% ===================================================================

%%K t10-apar-1 apar
\VUsaut{4.0mm}
\VUtitre{15.0pt}{60}{तक्ता क्रमांक १०}
\VUsaut{3.0mm}
\VUpk{11.4pt}{40}{समुद्र. --- बंदर.}
\VUsaut{3.0mm}

%%K t10-01-1 p
१. --- उन्हाळ्यात, काही जण डोंगरावर शांतता आणि गारवा शोधतात, तर दुसरे किनाऱ्याकडे जाणे पसंत
करतात. गेल्या वर्षी आम्ही तसेच केले, कारण आम्हाला \VUgras{महासागर} \textsuperscript{(1)} फार
आवडतो.

%%K t10-02-1 p
२. --- माझ्यापुरते सांगायचे तर, \VUgras{क्षितिजापर्यंत} \textsuperscript{(2)} पसरलेला पाण्याचा
तो अफाट विस्तार न्याहाळताना, आणि अमेरिकेला जाणारे प्रवासी ज्यांवर चढतात त्या अवाढव्य
\VUgras{आगबोटी} \textsuperscript{(3)} लांबच्या \VUgras{प्रवासाला} निघताना पाहताना, मला फार
आनंद होतो.

%%K t10-03-1 p
३. --- म्हणूनच माझा कल ओळखणारे माझे वडील सुट्टी घालवण्यासाठी नेहमी असा एखादा फार शांत किनारा
निवडतात, जो आपल्या मोठ्या \VUgras{लष्करी बंदरांपैकी} एकाजवळ असतो.

%%K t10-03-2 p
आमच्या गेल्या मुक्कामात \VUgras{आरमार} त्या अवाढव्य \VUgras{बंधाऱ्यामागे}
\textsuperscript{(4)} नांगरून आले, जो \VUgras{लष्करी बंदर} \textsuperscript{(5)} बंद करतो आणि
त्याचे रक्षण करतो; आणि मला मनासारख्या वेगवेगळ्या \VUgras{युद्धनौका} पाहता आल्या — पाण्याखाली
बुडून नाहीशी होणारी लहानशी \VUgras{पाणबुडी} \textsuperscript{(10)}, आणि ती अवाढव्य
\VUgras{चिलखती युद्धनौकाही} \textsuperscript{(9)}, जिला आजवर जवळजवळ केवळ
\VUgras{टॉर्पेडो नौकेच्या} \textsuperscript{(8)} हल्ल्यांचीच भीती होती.

%%K t10-tit-2 sub
\VUsaut{3.0mm}
\VUpk{10.2pt}{40}{लहान नौका.}
\VUsaut{3.0mm}

%%K t10-04-1 p
४. --- या शेवटच्या नौकेला जाड धातूच्या चिलखतातला कमजोर बिंदू मोठ्या \VUgras{कौशल्याने} शोधून
तिथे तो घातक \VUgras{टॉर्पेडो} डागणे आधीच जमत असे, ज्याच्या स्फोटाने नौका बुडते; तरीही ती
पाणबुडीपेक्षा कमी भीतिदायक होती, कारण टॉर्पेडोरोधक नौका तिचा सहज पाठलाग करतात.

%%K t10-05-1 p
५. --- मला चिलखत घातलेल्या वेगवान \VUgras{गस्तीनौकाही} \textsuperscript{(6)} पाहायला मिळाल्या
— खऱ्या चिलखती युद्धनौकेइतक्याच जवळजवळ अभेद्य —, कित्येक \VUgras{वाहतूक नौका}
\textsuperscript{(12)}, तीन-चार \VUgras{तोफनौका} \textsuperscript{(7)} आणि शेवटी एक
\VUgras{फ्रिगेट} \textsuperscript{(11)}.
\VUgras{नांगर उचलण्यासाठी ते नौकादल हालचाल करते तेव्हाचा देखावा सर्वात रोचक असतो.}

%%K t10-06-1 p
६. --- त्या पोलादाच्या भल्यामोठ्या ढिगांमध्ये आणि व्यापारी नौकादलात अजूनही वापरल्या जाणाऱ्या
डौलदार \VUgras{तीन डोलकाठ्यांच्या गलबतांमध्ये} किंवा रेखीव \VUgras{शिडाच्या गलबतांमध्ये}
\textsuperscript{(13)} केवढा बांधणीचा फरक! मी \VUgras{धक्क्यावर} \textsuperscript{(14)} फिरत
असताना ती सगळी शिडे फुलवून बंदरात शिरताना मी पाहिली. खरे तर \VUgras{वारा} उलटा असला
\VUgras{की} त्यांचे बरेच फायदे संपतात. बंदरात परत शिरण्यासाठी त्यांना सगळी शिडे उतरवावी
लागतात, आणि एका \VUgras{ओढनौकेची} \textsuperscript{(16)} गरज पडते, जी त्यांना आणून घाटाला
लावते — साध्या \VUgras{मालवाहू नौकांप्रमाणे} \textsuperscript{(17)}.

%%K t10-tit-3 sub
\VUsaut{3.0mm}
\VUpk{10.2pt}{40}{व्यापारी बंदर.}
\VUsaut{3.0mm}

%%K t10-07-1 p
७. --- मी ज्या बंदराबद्दल बोलतो त्याचे वैशिष्ट्य असे की त्यात केवळ अवाढव्य कामांनी कृत्रिम
रीतीने घडवलेले लष्करी बंदरच नाही, तर एका मोठ्या नदीच्या \VUgras{खाडीत} वसलेले एक अद्भुत
\VUgras{व्यापारी बंदरही} आहे.

%%K t10-08-1 p
८. --- दोन्ही बंदरे वेगळी करणारी ती जमिनीची चिंचोळी पट्टी तटबंदीने बांधली आहे आणि समुद्रात
पुढे आलेल्या \VUgras{तोफखान्यांच्या} आणि \VUgras{बुरुजांच्या} मालिकेने राखली जाते.
\VUgras{तटावर} \textsuperscript{(15)} फिरण्यासाठी खास परवाना लागतो.
\VUgras{जहाजबांधणी गोद्यांचे} \textsuperscript{(18)} अभियंता असलेल्या माझ्या काकांमुळे मला तो
सहज मिळाला, आणि विस्तीर्ण छपरांखाली बांधल्या जाणाऱ्या नौका पाहायला मी गेलो.

%%K t10-09-1 p
९. --- एका चिलखती युद्धनौकेचे \VUgras{जलावतरणही} मी पाहिले, आणि तो हेलावून टाकणारा क्षण मला
आठवतो, जेव्हा तो प्रचंड सांगाडा स्वतःच्याच भरवशावर सोडला जाऊन आपल्या पाळण्यासारख्या चौकटीवरून
घसरू लागला आणि नदीच्या पाण्यावर तरंगू लागला — तेव्हा सगळ्या गर्दीला जे वाटले ते.

%%K t10-10-1 p
१०. --- गोद्यांकडे जाण्यासाठी \VUgras{कवायतीचे मैदान} \textsuperscript{(19)} ओलांडावे लागते,
जिथे शिबंदीचे शिपाई रोज सराव करायला येतात, आणि लोखंडी \VUgras{पुलावरून} \textsuperscript{(20)}
जावे लागते, जो संचलनाचे मैदान \VUgras{शस्त्रागाराशी} \textsuperscript{(21)} जोडतो.

%%K t10-tit-4 sub
\VUsaut{3.0mm}
\VUpk{10.2pt}{40}{शस्त्रागार आणि गोद्या.}
\VUsaut{3.0mm}

%%K t10-11-1 p
११. --- काकांबरोबर मी \VUgras{तोफांनी} \textsuperscript{(22)}, \VUgras{तोफगोळ्यांनी}
\textsuperscript{(23)} आणि \VUgras{स्फोटक गोळ्यांनी} भरलेली ती मोठी इमारत पाहिली. तो
\VUgras{धातुकारखानाही} \textsuperscript{(24)} मी पाहिला, ज्यात आगबोटींचे \VUgras{बॉयलर}
\textsuperscript{(25)} तयार होतात.

%%K t10-12-1 p
१२. --- अगदी जवळ \VUgras{गोद्या} \textsuperscript{(26)} आहेत --- दुरुस्तीच्या गोद्या --- आणि
त्या फार लक्षणीय \VUgras{तरंगत्या गोद्या} \textsuperscript{(27)}, ज्यांत दुरुस्तीला आलेल्या
एका नौकेवर मला होडीचा \VUgras{पंखा} \textsuperscript{(28)}, \VUgras{तळकणा}
\textsuperscript{(29)} आणि \VUgras{सुकाणू} \textsuperscript{(30)} अगदी जवळून पाहता आले.

%%K t10-13-1 p
१३. --- व्यापारी बंदर लष्करी बंदराइतकेच अभ्यासण्याजोगे आहे, आणि नदी उलट चढून जाणाऱ्या
\VUgras{चाकांच्या नौका} \textsuperscript{(31)} जिथे बांधलेल्या असतात त्या घाटांवर मी कितीतरी
तास टक लावून घालवले — आणि \VUgras{डोलकाठीच्या याऱ्या} \textsuperscript{(32)} चालताना पाहत,
ज्या अवाढव्य ओझी पिसासारखी उचलतात.

%%K t10-tit-5 sub
\VUsaut{4.0mm}
\VUpk{10.2pt}{40}{वाटाड्या.}
\VUsaut{3.0mm}

%%K t10-14-1 p
१४. --- नांगरतळातून बाहेर पडणाऱ्या \VUgras{ट्रान्स-अटलांटिक नौकांचे} \textsuperscript{(34)}
मला कौतुक वाटले — त्यांची \VUgras{निशाणे} \textsuperscript{(35)} भरल्या हवेत फडकत, आणि एक कुशल
\VUgras{वाटाड्या} \textsuperscript{(36)} त्यांना चालवत, जो \VUgras{बोयांनी}
\textsuperscript{(40)} आणि खुणादर्शकांनी आखलेल्या \VUgras{कालव्यातून} अद्भुत कौशल्याने नेतो,
आपल्या एकाच चौकोनी शिडाने कष्टाने वळवल्या जाणाऱ्या \VUgras{सपाट तळाच्या नौका}
\textsuperscript{(41)} चुकवत. \VUgras{पुढच्या भागावर} \textsuperscript{(37)} --- नाळेवर ---
मला दुसऱ्या वर्गाचे \VUgras{कक्ष} \textsuperscript{(39)} ओळखू आले, आणि \VUgras{मागच्या भागावर}
\textsuperscript{(38)} --- पिछाडीवर --- पहिल्या वर्गाच्या प्रवाशांचे दिवाणखाने.

%%K t10-15-1 p
१५. --- कधी कधी मी \VUgras{मालवाहू नौकेची} \textsuperscript{(42)} भरणीही पाहिली, जिच्यात
\VUgras{गुदामातला} \textsuperscript{(43)} आणि \VUgras{सागरी स्थानकातला} \textsuperscript{(44)}
माल नुकताच रचला गेला होता — \VUgras{घाटावरच्या लोहमार्गाच्या} \textsuperscript{(45)} असंख्य
गाड्यांनी आणलेला.

%%K t10-tit-6 sub
\VUsaut{3.0mm}
\VUpk{10.2pt}{40}{वल्हवण्याची मौज.}
\VUsaut{3.0mm}

%%K t10-16-1 p
१६. --- मी अनेकदा \VUgras{तरंगत्या गोदीत} \textsuperscript{(53)} होडी वल्हवली, जिच्यात
\VUgras{होड्या} \textsuperscript{(46)} आश्रयाला येतात, आणि \VUgras{वल्हे}
\textsuperscript{(47)} चालवणे माझ्यासाठी आनंदाचे होते. मी \VUgras{डेकवर}
\textsuperscript{(48)} चढलो — एका \VUgras{शिडाच्या होडीच्या} \textsuperscript{(49)} —,
\VUgras{दोरखंडांनी} \textsuperscript{(52)} \VUgras{डोलकाठीवर} \textsuperscript{(51)}
\VUgras{आडदांड्यांपर्यंत} \textsuperscript{(50)} चढलो, आणि मग \VUgras{लहान होडीतून}
\textsuperscript{(54)} माझा फेरफटका पुन्हा केला — जड \VUgras{मासेमारीच्या होड्यांमधून}
\textsuperscript{(55)}, जिथे \VUgras{जाळी} \textsuperscript{(56)} वाळत असतात.

%%K t10-17-1 p
१७. --- \VUgras{घाटावर} \textsuperscript{(58)} माझ्यासमोरून अत्यंत निरनिराळा \VUgras{माल}
\textsuperscript{(59)} सरकत गेला, \VUgras{पेट्यांत} \textsuperscript{(60)} किंवा
\VUgras{गासड्यांत} \textsuperscript{(61)} भरलेला. मोठ्या होड्यांचे \VUgras{ओझे}
\textsuperscript{(63)} तोच होता, आणि जबरदस्त \VUgras{याऱ्या} \textsuperscript{(62)} तो उचलून
\VUgras{ट्रकांवर} \textsuperscript{(64)} ठेवत होत्या, तर \VUgras{मालवाहतूकदार}
\VUgras{जकात कार्यालयात} \textsuperscript{(65)} त्याची नोंद देत होते आणि कर भरत होते.

%%K t10-18-1 p
१८. --- शेवटी, मी \VUgras{फिरत्या पुलाशी} \textsuperscript{(66)} किंवा \VUgras{जलद्वाराशी}
\textsuperscript{(69)} पोहोचलो, \VUgras{गाळ काढणाऱ्या यंत्रासमोरून} \textsuperscript{(68)} जात
— जे \VUgras{कालवा} \textsuperscript{(67)} साफ करते —, आणि \VUgras{संकेतस्तंभाजवळ}
\textsuperscript{(70)} धक्क्यावर उतरलो.

%%K t10-19-1 p
१९. --- त्या धक्क्याच्या दुसऱ्या बाजूलाच त्या \VUgras{नावा} \textsuperscript{(71)} लागतात,
ज्या आरमारातले अधिकारी जमिनीवर आणतात. खलाशी तालात आपली वल्ही उचलताना पाहणे हा कौतुकास्पद
देखावा असतो, आणि तेवढ्यात \VUgras{लाटेवर} \textsuperscript{(72)} तरंगणारी होडी उतरण्याच्या
जिन्याला सहज लागते. या जागी \VUgras{भरती-ओहोटी} आणि \VUgras{किनाऱ्यावरच्या}
\textsuperscript{(73)} \VUgras{भरतीची} व \VUgras{ओहोटीची} हालचाल अधिक चांगली पाहता येते.

%%K t10-tit-7 sub
\VUsaut{3.0mm}
\VUpk{10.2pt}{40}{क्लब.}
\VUsaut{3.0mm}

%%K t10-20-1 p
२०. --- घरी परतण्यासाठी मी \VUgras{विजेची ट्राम} \textsuperscript{(74)} वापरली, जिचा
\VUgras{थांबा} \textsuperscript{(75)} अधिकाऱ्यांच्या क्लबासमोर उभ्या केलेल्या
\VUgras{खांबाजवळ} आहे.

%%K t10-20-2 p
क्लबाच्या \VUgras{सज्ज्यावरून} \textsuperscript{(76)} — \VUgras{नांगरांनी}
\textsuperscript{(77)}, तोफांनी आणि गोळ्यांनी सजवलेल्या — मी अनेकदा नौदल अधिकाऱ्यांचा भव्य
मेळावा पाहिला : \VUgras{लेफ्टनंट} \textsuperscript{(78)} आपल्या तलवारींसह,
\VUgras{तोफखान्याचे कप्तान} \textsuperscript{(79)} आणि \VUgras{पायदळाचे}
\textsuperscript{(80)} कप्तान आपल्या \VUgras{समशेरींसह}. त्यांच्यामागे एक \VUgras{खलाशी}
\textsuperscript{(81)} असे, जो लांबचे काही ओळखायचे असले की त्यांना \VUgras{दुर्बीण} आणून देई.

%%K t10-21-1 p
२१. --- तरुण \VUgras{उपलेफ्टनंट} \textsuperscript{(82)} आपल्या \VUgras{सेनापतीकडून}
\textsuperscript{(83)} किंवा \VUgras{सरखेलाकडून} \textsuperscript{(84)} सूचना घेतानाही मी
पाहिले, आणि तेवढ्यात एक \VUgras{तांडेल} \textsuperscript{(85)} त्या नौकेवर नेण्यासाठी वाट पाहत
उभा असे.

%%K t10-22-1 p
२२. --- त्या सज्ज्यावरून दृश्य भव्य आहे : आपल्या \VUgras{तंबूंसह} \textsuperscript{(86)} आणि
\VUgras{कपडे बदलायच्या खोल्यांसह} \textsuperscript{(87)} सगळा किनारा नजरेखाली येतो, आणि
आंघोळीच्या वेळी \VUgras{स्नान करणारे} \textsuperscript{(88)} दिसतात, \VUgras{कॅसिनोसमोर}
\textsuperscript{(89)} आनंदाने दंगा करणारे.

%%K t10-23-1 p
२३. --- किनाऱ्याच्या टोकाला जमीन एक लहानसे \VUgras{खडकाळ} \VUgras{द्वीपकल्प}
\textsuperscript{(91)} तयार करते, जिथे \VUgras{मोठ्या लाटा} \textsuperscript{(92)} येऊन फुटतात
आणि ज्यावर एक \VUgras{दीपगृह} \textsuperscript{(93)} बांधले आहे, जे रात्री नावाड्यांना वाट
दाखवते. ते द्वीपकल्प जमिनीशी केवळ एका फार अरुंद \VUgras{संयोगभूमीने} \textsuperscript{(90)}
जोडलेले आहे.

%%K t10-tit-8 sub
\VUsaut{3.0mm}
\VUpk{10.2pt}{40}{किनाऱ्यांचे सर्वसाधारण दृश्य.}
\VUsaut{3.0mm}

%%K t10-24-1 p
२४. --- \VUgras{समुद्रकडा} \textsuperscript{(94)} लांबवर पसरतो, समुद्राने चिरलेला, जो
त्याच्यात लहरीने \VUgras{उपसागरांची} \textsuperscript{(95)} आणि \VUgras{भूशिरांची} मालिका
कोरतो आणि त्याला अत्यंत चित्रमय करतो.

%%K t10-24-2 p
\VUgras{पठारावर} \textsuperscript{(96)}, जे \VUgras{सामुद्रधुनीवर} \textsuperscript{(98)} नजर
ठेवते, एक फार महत्त्वाचा \VUgras{किल्ला} \textsuperscript{(97)} आणि काही « बंगले » आहेत.

%%K t10-25-1 p
२५. --- ती सामुद्रधुनी खंडापासून एक फार धोकादायक \VUgras{बेट} \textsuperscript{(99)} वेगळे
करते, जे \VUgras{खडकांनी} \textsuperscript{(100)} आणि \VUgras{फुटणाऱ्या लाटांनी} वेढलेले आहे,
जिथे कधी कधी होड्या आणि मोठ्या नौकाही फुटतात. मदत तत्पर असूनही आणि \VUgras{बचावनौका} भराभर
पोहोचत असूनही, ते \VUgras{जहाजभंग} भयंकर असतात, आणि संपातकाळच्या \VUgras{वादळांत} कित्येक नौका
माणसांसह आणि मालासह बुडाल्या.

%%K t10-25-2 p
हे शेजारपण असूनही \VUgras{ते बंदर} खलाशांची फार वर्दळ असलेले आहे.

\VUsaut{6.0mm}
\VUfilet{22mm}
